\begin{tikzpicture}[
  >=Stealth,
  node distance=1.5cm and 2cm,
  data/.style={draw, fill=blue!10, minimum width=2.4cm, minimum height=1.4cm, align=center, font=\small\bfseries, rounded corners=2pt},
  module/.style={draw, rounded corners=4pt, fill=blue!20, minimum width=3.2cm, minimum height=1.6cm, align=center, font=\small\bfseries},
  core/.style={draw, very thick, rounded corners=2pt, fill=orange!30, minimum width=3.4cm, minimum height=1.8cm, align=center, font=\small\bfseries, double=orange!80, double distance=1.5pt},
  output/.style={draw, rounded corners=4pt, fill=green!20, minimum width=2.4cm, minimum height=1.4cm, align=center, font=\small\bfseries},
  arrow/.style={-{Stealth[length=2.5mm, width=2mm]}, thick, color=black!80},
  maskarrow/.style={-{Stealth[length=2.5mm, width=2mm]}, thick, color=red!60!black, dashed},
  label/.style={font=\scriptsize\itshape, text=black!70},
  masklabel/.style={font=\scriptsize\itshape, text=red!60!black},
]

% ================= 节点定义 =================
% 输入节点
\node[data, fill=gray!20] (input) at (0,0) {Light Field Input\\ \scriptsize $(B, 9, 9, H, W, 3)$};

% 上分支：空间特征提取
\node[module] (epi_ext) at (3.8, 2.2) {LF Input \& EPI\\Extraction\\ \scriptsize 4-Dir EPI};

% 下分支：角度频域分析（核心创新）
\node[core] (ang_freq) at (3.8, -2.2) {Angular Frequency\\Analysis\\ \scriptsize (AngularFreqNet)};

% 下分支：双掩码生成（核心创新）
\node[core] (dual_mask) at (8.2, -2.2) {Dual-Mask\\Generation\\ \scriptsize $M_{mat}, M_{conf}$};

% 上分支：深度回归（核心创新）
\node[core] (epinet) at (8.2, 2.2) {EPINet4Dir Depth\\Regression\\ \scriptsize $D_{init}$};

% 汇合：统一物理模型深度优化（核心创新）
\node[core, minimum width=3.6cm] (refine) at (12.5, 0) {Unified Physical\\Depth Refinement};

% 输出节点
\node[output] (output) at (16.2, 0) {Final Depth Map\\ \scriptsize $(B, 1, H, W)$};

% ================= 数据流连接 =================
% 1. 输入分发
\draw[arrow] (input.east) -- ++(1.2,0) coordinate (split);
\draw[arrow] (split) |- (epi_ext.west) node[pos=0.25, above, label] {Spatial};
\draw[arrow] (split) |- (ang_freq.west) node[pos=0.25, below, label] {Angular};
\fill (split) circle (1.5pt);

% 2. 上分支前向传播
\draw[arrow] (epi_ext.east) -- (epinet.west) node[midway, above, label] {4-Dir EPI};

% 3. 下分支前向传播
\draw[arrow] (ang_freq.east) -- (dual_mask.west) node[midway, above, label] {$F_{low}, F_{high}$};

% 4. 掩码控制流（双掩码 -> 深度回归）
\draw[maskarrow] (dual_mask.north) -- (epinet.south) node[midway, right, masklabel] {$M_{conf}$};

% 5. 汇聚到物理优化模块
\draw[arrow] (epinet.east) -- ++(1.2,0) |- (refine.north) node[pos=0.25, above, label] {$D_{init}$};
\draw[maskarrow] (dual_mask.east) -- ++(1.8,0) |- (refine.south) node[pos=0.25, below, masklabel] {$M_{mat}$};

% 6. 最终输出
\draw[arrow] (refine.east) -- (output.west) node[midway, above, label] {$D_{final}$};

% ================= 关键公式标注 =================
\node[font=\tiny\itshape, text=black!60, align=center] at (3.8, -3.6) {$F(\omega) = \sum I(\theta) e^{-j \frac{2\pi}{N} \omega \theta}$};
\node[font=\tiny\itshape, text=black!60, align=center] at (8.2, -3.6) {$M = \sigma(\text{Conv}(F_{low} \oplus F_{high}))$};
\node[font=\tiny\itshape, text=black!60, align=center] at (12.5, -1.6) {$D_{final} = D_{init} + \lambda \nabla \cdot (M_{mat} \odot \nabla D_{init})$};

\end{tikzpicture}